The rapid advancement of large language models (LLMs) has transformed scientific discovery through agentic systems that autonomously conduct experiments and test hypotheses~\citep{lu2024ai, yamada2025ai, novikov2025alphaevolve, zhang2025darwin}. These frameworks leverage LLMs as sophisticated mutation operators, iteratively refining candidate solutions with successful variants propagating through successive generations. This methodology has proven effective across domains such as competitive programming \citep{li2022competition}, mathematical optimization \citep{romera2024mathematical}, and automated agentic design \citep{hu2024automated}.
%
However, current implementations face significant practical limitations. The primary challenge is substantial sample inefficiency as existing approaches typically require thousands of evaluations, making them computationally expensive and time-consuming. This inefficiency stems from naive exploration strategies that fail to effectively leverage accumulated knowledge from previous generations. Additionally, most leading systems remain closed-source, creating barriers to reproducibility and limiting community-driven improvements.
%
\ouralgo addresses these challenges through three key algorithmic innovations that work synergistically to enhance sample efficiency. Our adaptive parent and LLM sampling intelligently balances exploration of novel regions with exploitation of known high-quality areas. Next, our code proposal novelty rejection sampling ensures efficient program mutations. Finally, our bandit-based LLM ensemble selection strategy dynamically adapts to the evolving state of the sampled archive parents and inspiration programs.
%
Experimental validation across diverse domains demonstrates substantial improvements in both efficiency and solution quality, with \ouralgo achieving state-of-the-art results using orders of magnitude fewer evaluations than existing approaches. By releasing our complete implementation as open-source software, we aim to democratize access to advanced evolutionary discovery tools and enable broad community contributions. In summary:

\begin{enumerate}
    \item We introduce \ouralgo, an evolutionary framework that substantially improves sample efficiency through three key algorithmic innovations: a novel parent program sampling strategy, code novelty rejection-sampling, and adaptive performance-based LLM ensemble selection.
    \item We demonstrate \ouralgo's ability to innovate beyond human and LLM-generated solutions with comprehensive experimental validation across four distinct problem domains: mathematical optimization (circle packing), agentic design (AIME tasks), competitive programming (ALE-Bench), and LLM training design (mixture-of-expert load balancing loss). 
    %\rob{add numbers on results?} EDO: I think it could be too early to add concrete numbers without properly contextualizing what these numbers represent - but feel free to modify the last sentence if you think otherwise ;)
    \item We release \ouralgo as open-source software under the Apache 2.0 license, including implementation details, and an interactive visualization tool for monitoring the search process.
\end{enumerate}
